% GENERATED FILE. Edit canonical lessons, then run npm run generate:books.
% canonical-source-hash: fnv1a64:4d3eb6a4d2a5872d
% canonical-lessons: FR-C137-commencer, FR-C137-visiter, FR-C137-repondre, FR-C137-mettre, FR-C137-croire

\chapter{To begin, To visit (a place), To answer, To put, To believe}
\label{ch:fr-to-begin-to-visit-a-place-to-answer-to-put-to-believe}
\hlchaptermodality{137}

\begin{chapteropening}
\textbf{By the end of this chapter:} \emph{I can say to begin, to visit (a place), to answer, to put, to believe.}

Complete the last lesson of chapter 137: I can say to begin, to visit (a place), to answer, to put, to believe.
\end{chapteropening}

\section[commencer]{commencer — to begin}

\label{lesson:FR-C137-commencer}



\begin{quote}
Before the new one: say the French for ninety, then the French for a thousand.
\end{quote}

\subsection*{You'll want to know: commencer}
\textbf{commencer} — \textquotedblleft{}to begin\textquotedblright{}.

\begin{cousinweb}
\textbf{commencer} $\leftarrow$ Old French \emph{comencier} $\leftarrow$ Vulgar Latin \emph{*cominitiare}, from \emph{com-} and Latin \emph{initiare}, \textquotedblleft{}to begin\textquotedblright{} — inherited. (EtymDB 2.1)
\end{cousinweb}

\begin{grammarlens}[title={What you've built}]
One more word for this chapter.
\end{grammarlens}

\subsection*{Your turn}
\begin{itemize}
\raggedright
  \item \emph{Say it:} \emph{commencer}
  \item \emph{Say it:} \emph{commencer}, once more
  \item \emph{Say it:} say \emph{commencer}
  \item \emph{Recall:} say the French for ninety, then the French for a thousand, then say \emph{commencer} again
\end{itemize}

\begin{tcolorbox}[breakable,colback=teal!4,colframe=teal!35!black,title={Before you move on}]
What does commencer mean? (\textquotedblleft{}To begin\textquotedblright{}.) Say it once more.
\end{tcolorbox}
\section[visiter]{visiter — to visit (a place)}

\label{lesson:FR-C137-visiter}



\begin{quote}
Before the new one: say the French for a thousand, then the French for to begin.
\end{quote}

\subsection*{You'll want to know: visiter}
\textbf{visiter} — \textquotedblleft{}to visit (a place)\textquotedblright{}.

\begin{cousinweb}
\textbf{visiter} $\leftarrow$ Latin \emph{visitare}, \textquotedblleft{}to go to see\textquotedblright{}, from \emph{videre}, \textquotedblleft{}to see\textquotedblright{} — a learned borrowing. (EtymDB 2.1)
\end{cousinweb}

\begin{grammarlens}[title={What you've built}]
One more word for this chapter.
\end{grammarlens}

\subsection*{Your turn}
\begin{itemize}
\raggedright
  \item \emph{Say it:} \emph{visiter}
  \item \emph{Say it:} \emph{visiter}, once more
  \item \emph{Say it:} say \emph{visiter}
  \item \emph{Recall:} say the French for a thousand, then the French for to begin, then say \emph{visiter} again
\end{itemize}

\begin{tcolorbox}[breakable,colback=teal!4,colframe=teal!35!black,title={Before you move on}]
What does visiter mean? (\textquotedblleft{}To visit (a place)\textquotedblright{}.) Say it once more.
\end{tcolorbox}
\section[répondre]{répondre — to answer}

\label{lesson:FR-C137-repondre}



\begin{quote}
Before the new one: say the French for to begin, then the French for to visit.
\end{quote}

\subsection*{You'll want to know: répondre}
\textbf{répondre} — \textquotedblleft{}to answer\textquotedblright{}.

\begin{cousinweb}
\textbf{répondre} $\leftarrow$ Old French \emph{respondre} $\leftarrow$ Latin \emph{respondere}, \textquotedblleft{}to answer, to pledge back\textquotedblright{} — inherited. (EtymDB 2.1)
\end{cousinweb}

\begin{grammarlens}[title={What you've built}]
One more word for this chapter.
\end{grammarlens}

\subsection*{Your turn}
\begin{itemize}
\raggedright
  \item \emph{Say it:} \emph{répondre}
  \item \emph{Say it:} \emph{répondre}, once more
  \item \emph{Say it:} say \emph{répondre}
  \item \emph{Recall:} say the French for to begin, then the French for to visit, then say \emph{répondre} again
\end{itemize}

\begin{tcolorbox}[breakable,colback=teal!4,colframe=teal!35!black,title={Before you move on}]
What does répondre mean? (\textquotedblleft{}To answer\textquotedblright{}.) Say it once more.
\end{tcolorbox}
\section[mettre]{mettre — to put; to put on}

\label{lesson:FR-C137-mettre}



\begin{quote}
Before the new one: say the French for to visit, then the French for to answer.
\end{quote}

\subsection*{You'll want to know: mettre}
\textbf{mettre} — \textquotedblleft{}to put; to put on\textquotedblright{}.

\begin{cousinweb}
\textbf{mettre} $\leftarrow$ Old French \emph{metre} $\leftarrow$ Latin \emph{mittere}, \textquotedblleft{}to send, to let go\textquotedblright{} — inherited. (EtymDB 2.1)
\end{cousinweb}

\begin{grammarlens}[title={What you've built}]
One more word for this chapter.
\end{grammarlens}

\subsection*{Your turn}
\begin{itemize}
\raggedright
  \item \emph{Say it:} \emph{mettre}
  \item \emph{Say it:} \emph{mettre}, once more
  \item \emph{Say it:} say \emph{mettre}
  \item \emph{Recall:} say the French for to visit, then the French for to answer, then say \emph{mettre} again
\end{itemize}

\begin{tcolorbox}[breakable,colback=teal!4,colframe=teal!35!black,title={Before you move on}]
What does mettre mean? (\textquotedblleft{}To put; to put on\textquotedblright{}.) Say it once more.
\end{tcolorbox}
\section[croire]{croire — to believe}

\label{lesson:FR-C137-croire}



\begin{quote}
Before the new one: say the French for to answer, then the French for to put; to put on.
\end{quote}

\subsection*{You'll want to know: croire}
\textbf{croire} — \textquotedblleft{}to believe\textquotedblright{}.

\begin{cousinweb}
\textbf{croire} $\leftarrow$ Old French \emph{creire} (EtymDB 2.1) $\leftarrow$ Latin \emph{credere}, \textquotedblleft{}to believe\textquotedblright{} (Lewis \& Short) — inherited.
\end{cousinweb}

\begin{grammarlens}[title={What you've built}]
One more word for this chapter.
\end{grammarlens}

\subsection*{Your turn}
\begin{itemize}
\raggedright
  \item \emph{Say it:} \emph{croire}
  \item \emph{Say it:} \emph{croire}, once more
  \item \emph{Say it:} say \emph{croire}
  \item \emph{Recall:} say the French for to answer, then the French for to put; to put on, then say \emph{croire} again
\end{itemize}

\begin{tcolorbox}[breakable,colback=teal!4,colframe=teal!35!black,title={Before you move on}]
What does croire mean? (\textquotedblleft{}To believe\textquotedblright{}.) Say it once more.
\end{tcolorbox}
